% Seminar information shared by 01-... and 02-...
\newcommand{\includetitle}{Generative modeling of microstructures: From image to texture to graphs}
\newcommand{\includeauthor}{Gowtham Nimmal Haribabu}
\newcommand{\includeaffiliation}{TU Delft}
\newcommand{\includedate}{October 15, 2026}
\newcommand{\includetime}{2 -- 3 pm}
\newcommand{\includevenue}{Amphithéâtre Pole Mecanique}
\newcommand{\includephoto}{images/nh}
\newcommand{\includeabstract}{The mechanical response of engineering materials emerges from their underlying microstructure, and their evolution during processing and deformation. However, establishing the process-microstructure-property relationships remains challenging because microstructures are high-dimensional and stochastic. Conventional simulation and statistical approaches provide valuable insights, but can become computationally expensive or difficult to use for rapid forward prediction and inverse design.
Generative machine learning offers a complementary route by learning probability distributions over microstructures and enabling conditional generation directly from processing variables or target properties. In this talk, I will present three generative approaches that operate on different microstructural representations. First, I will discuss a conditional diffusion model for microstructure images. Then I will present TexFlow, a conditional normalizing flow framework for crystallographic texture of grain orientations while respecting the geometry of SO(3). Finally I will introduce graph-based modeling of three-dimensional microstructures. 
I will close by highlighting the key opportunities and remaining challenges in developing generative models that are data-efficient, physically meaningful and useful for predictive and inverse materials design. 
}
\newcommand{\includebio}{Gowtham received a B. Tech in Metallurgical and Materials Engineering from the National Institute of Technology, Tiruchirappalli, a M. Tech in Metallurgical and Materials Engineering from the Indian Institute of Technology, Kharagpur, and a Ph. D in Materials from the Indian Institute of Science, Bengaluru. He then did his postdoc in Residual stress measurement at University of Manchester and currently a research associate at Materials Science and Engineering, TU Delft. His research interest includes applied machine learning for materials, biomechanics, additive manufacturing. }
